\subsection{Overview of the models}

As seen in the previous section \ref{methodology}, the model of the community is built in a modular way. We first define the model of the different devices and flexibility services, then the member model and finally the community model. So in this part, we will present in detail how the devices are separated in different categories and how to define their model. Then we will see exactly how the aggregated model are built and define the different objective functions. This model is multi-objective as we want to optimize the community based on sociological profiles of the members. Therefore, we need to build different objective functions, and we need a way to compare them in order to plot the Pareto front.

\subsection{Flexibility services model}

We already identified the different type of flexibility services in the state of the art \ref{section_optimization}. In this work, we will focus on the following ones: the white goods (with a fixed cycle that can be shifted in time), the thermal load (with a comfort range) and the battery storage and electric vehicle (with a storage capacity and a maximum charge and discharge power), and finally the fixed loads with no flexibility.

For each device, what is important is to define one variable representing the consumption profile, that will have the same shape for all the devices, so that they can be summed directly. This variable will be noted $P_{device}$, define as a set of variable for each time step $t \in T$.

\subsubsection{Fixed loads}

Here, the profile is simple enough, the consumption profile is given as a parameter, so the only constraint is  :

\begin{equation} 
    P_{device}[t] = P_{fixed}[t], \forall t \in T
\end{equation}
With $P_{fixed}$ the fixed consumption profile of the device.

\subsubsection{Thermal load}

For defining the thermal load, we have a consumption profile in the form of a range between a minimum and a maximum consumption. We define the maximum consumption as the optimal comfort one. We will then define a variable that compute the gap between the optimal comfort and the chosen power consumption, that will be noted $p_{comfort}$. So we have two constraints defined as follows:
\begin{align} 
    P_{min}[t] \leq P_{device}[t] \leq P_{max}[t], \forall t \in T \\
    p_{comfort}[t] = P_{max}[t] - P_{device}[t], \forall t \in T
\end{align}

\subsubsection{White goods}

For the white goods, we have a fixed cycle that can be shifted in time. So we define a variable 
$x_{start}$ that will be equal to 1 if the cycle starts at time step $t$ and 0 otherwise. This time, we define a comfort time start as well as a time range in which the cycle needs to be executed. Using this, the gap between the starting time and the time wanted ($t_{wanted}$) is defined as a time comfort variable ($t_{comfort}$). In order to use a MILP (mixed integer linear programming) model, we defined two variables in order to define 
$t_{comfort}$, $t_{plus}$ and $t_{minus}$ for positive and negative gap values. 
A parameter $P_{cycle}$ that will be equal to the consumption profile of the cycle is given. For doing so, we define a new time set $T_{cycle}$ that will be all the possible time step of the cycle. The comfort constraints are then defined as follows:


\begin{align}
    &\sum_{t \in T_{cycle}} x_{start}[t] = 1 \\
    &\sum_{t \in T_{cycle}} (t+1) * x_{start}[t] - t_{wanted} \leq t_{plus} \\
    &t_{wanted} - \sum_{t \in T_{cycle}} (t+1) * x_{start}[t] \leq t_{minus} \\
    &t_{comfort} = t_{plus} + t_{minus} 
\end{align}
These constraints ensure that the cycle is executed once and that the comfort time variable is equal to the gap between the wanted time and the actual start time of the cycle.


Then the power consumption constraints are defined as follows:
\begin{align}
    &S_{power}[t] = [max(t-cycle, t_{wanted}-t_{range}), min(t, t_{wanted}+t_{range}-cycle)] \\
    &P_{device}[t] = \sum_{p \in S_{power}[t]} P_{cycle}[t-p] * x_{start}[p], \forall t \in T \\
    &P_{device}[t] = 0, \forall t \notin T_{cycle} 
\end{align}
With $cycle$ being the duration of the cycle. These power constraints ensure that the power consumption of the device is equal to the cycle profile when the cycle is running and 0 otherwise.

\subsubsection{Battery storage and electric vehicle}

Finally, comes the battery storage and electric vehicle. The only difference between the two is that we cannot operate the electric vehicle when it is not at home. Also, a minimum state of charge is required before leaving and a new value is defined after coming back home. The choice has been to define a battery as a fixed electric vehicle, so both of them have the same constraints. Here are the constraints for the battery storage and electric vehicle:

        
\begin{align}
    &E[t] = E[t-1] + P_{plus}[t] * \eta_{charge} - P_{minus}[t] / \eta_{dcharge}, \forall t \in T_{Home} \\
    &P_{device}[t] = P_{plus}[t] - P_{minus}[t], \forall t \in T_{Home} \\
    &P_{device}[t] = 0, \forall t \in T_{NotHome} \\
    &E[0] = E_0 \\
    &E[end] = E_{end} \\
    &E[t] \geq E_{min}[t], \forall t \in T_{leave} \\
    &E[t] = E[t'] - \Delta E, \forall t \in T_{back} \text{ and } t' \text{equal to the corresponding leave time}
\end{align}
With $E$ the state of charge of the battery, $P_{plus}$ and $P_{minus}$ the charge and discharge power. They are here to differentiate between the charging and discharging phase. Those values are positive. $\eta_{charge}$ and $\eta_{dcharge}$ the charge and discharge efficiency, $E_0$ the initial state of charge, $E_{end}$ the final state of charge. $E_{min}$ is the list of the minimum state of charge at the different leave time, $\Delta E$ is the energy needed for the trip and $T_{leave}$ and $T_{back}$ are the time steps of leaving and coming back home. Finally, $T_{Home}$ and $T_{NotHome}$ are the time steps when the electric vehicle is at home or not respectively. 

\subsection{Prosumer model}

Each member of the community is defined as a prosumer, meaning that they can both produce and consume energy. Their consumption profile is defined by the different devices they have, while their production profile is given as a parameter. This model supposes that we know beforehand the production profile of each member, by a predicting model for example. We then need to add to the model the power exchange with the community and with the grid.

The first step here is to define the link between the different devices and the total consumption profile of the member. This is done by defining a variable $P_{member}$ and $P_{battery}$ this way : 

\begin{align}
    &P_{member}[t] = \sum_{d \in D} P_{d}[t], \forall t \in T \\ 
    &P_{battery}[t] = \sum_{d \in D_{battery}} P_{d}[t], \forall t \in T
\end{align}
With $D$ the set of devices of the member other than the battery storage and electric vehicle, and $D_{battery}$ the set of battery storage and electric vehicle devices.

For the power exchange with the community, we will differentiate between the power coming from the community and the power sent to the community. Also, we enable exchanges within the community only from the produced energy, coming either from the power source or from the battery storage. For doing so, again we need to define a positive and negative variable for the battery power ($P_{bat+}$ and $P_{bat-}$). So we define a variable $P_{shared}$ and a variable $P_{self}$ for the shared and self-consumed power respectively. The constraints for sharing power are defined as follows:

\begin{equation}
    P_{shared}[t] + P_{self}[t] = P_{prod}[t] + P_{bat-}[t], \forall t \in T
\end{equation}

At the member scale, we will not consider in the power balance the power sent to the community, as it will be done at the community level. So for the power balance we consider only the power coming from the community defined as $P_{received}$. For the power exchanged with the grid, we want to differentiate between the power bought from the grid and the power sold to the grid for the objective function definition. So we define two variables $P_{grid+}$ and $P_{grid-}$ for the power bought and sold to the grid respectively. The power balance constraint is defined as follows:
\begin{equation}
    P_{grid+}[t] - P_{grid-}[t] = P_{member}[t] + P_{bat+}[t] - P_{received}[t] - P_{self}[t], \forall t \in T
\end{equation}
We can also sum the comfort variables of the different devices to have a global comfort variable for the member.

By adding the objective functions to this model, we can complete the fully decentralized optimization model. These objective functions will be defined just after in \ref{objective_function}.

\subsection{Community model}
The community model is built by aggregating the different member models and adding the constraints related to the power exchange between the members. For doing so, we need to define a variable $P_{exchange}[i, j, t]$ for each time step $t$ and each pair of members $i$ and $j$. It represents the power sent from member $i$ to member $j$ at time step $t$ and is positive. The power balance constraints for the community are then defined as follows: 

\begin{align}
    &P_{received, m}[t] = \sum_{j \in M} P_{exchange}[j, m, t], \forall m \in M, \forall t \in T \\
    &P_{exchange}[i, i, t] = 0, \forall i \in M, \forall t \in T 
\end{align}
With $M$ the set of members of the community. The power consumed, produced as well as the comfort variables can then be summed directly from the member models. By adding an objective function to this model, we can optimize the system in a fully centralized way.

\subsection{Objective function definition}
\label{objective_function}

In \cite{these_sorel}, they use a vector of the sociological profile of each member of the community in order to define a multi-objective optimization problem as explained in the state of the art \ref{section_sociological_profiles}. In this work, we reproduced this idea. They used a vector of 3 dimensions, the economic dimension, the environmental dimension and the independence dimension.

The way the model is built here, the independence dimension is linked very strongly to the other two dimensions, that already try to maximize the self-consumption of the members. Furthermore, the flexibility could very well impact the comfort of the members. So the idea has been to replace the independence dimension by a comfort dimension. The independence dimension could be more interesting in a model as defined in \cite{article_5} where the community become disconnected from the main grid.

For the economic dimension, it is easy to define the objective function as the cost of the energy bought from the grid minus the revenue from the energy sold to the grid. For the environmental dimension, we can associate to the grid energy a carbon footprint and another one to the energy consumed by the devices. For example, in France, the carbon footprint of the grid energy is in average equal to 40 gCO2/kWh in 2022 while for the same year, in France, the carbon footprint of the energy produced by a PV panel is around 25 gCO2/kWh. Of course a more accurate study could be done by looking exactly at the installation of each member of the community. The main point is that the two objective functions defined like this will remain linear and that we could imagine some coefficient to be able to compare the two dimensions easily. Here are the equations of such objective functions :

\begin{align}
    ElectricityPrice = \sum_{t \in T} &\big(P_{grid+}[t] \cdot BuyPrice[t] - P_{grid-}[t] \cdot SellPrice[t] 
        \nonumber \\&+ P_{received} \cdot CommunityPrice\big)  \\ \nonumber \\
    CarbonFootprint = \sum_{t \in T} &\big(P_{grid+}[t] \cdot GridFootprint[t] \nonumber \\
    &+ P_{self}[t] \cdot SelfFootprint 
        \nonumber\\
        &+ P_{shared}[t] \cdot SharedFootprint\big)
\end{align}
With $BuyPrice$ and $SellPrice$ the price of buying and selling energy to the grid, $CommunityPrice$ the price of energy received from the community, $GridFootprint$, $SelfFootprint$, and $SharedFootprint$ the carbon footprint of the grid energy, the self-consumed energy, and the shared energy respectively. For simplicity's sake, the prices and carbon footprint have been considered as constant over time. In this study we took $CommunityPrice$ equal to 0. But this number could be used to study the effect of dynamic pricing within the community for example.


But for the comfort, we have the choice about how to build the objective function, and it seems to be harder to compare with the other. The choice made in this work is to make a linear objective function, and then we will compute before each optimization references value of the community (or the member) for each dimension. We will then divide the value of the objective function by this reference value in order to pass in per unit values and be able to compare the objective functions. So here is the equation of the comfort objective function:

\begin{equation}
    Comfort = \sum_{t \in T} p_{comfort}[t] \cdot coef_p + t_{comfort} \cdot coef_t \cdot |T|
\end{equation}
With $coef_p$ and $coef_t$ the coefficient for the power comfort and time comfort respectively. The choice of these coefficients will determine the importance of each term in the comfort objective function. For this study they have been set to 1 each.

Finally the complete objective function is defined as follows:

\begin{equation}
    Objective = w_{eco} \cdot \frac{ElectricityPrice}{ref_{eco}} + w_{env} \cdot \frac{CarbonFootprint}{ref_{env}} + w_{comfort} \cdot \frac{Comfort}{ref_{comfort}}
\end{equation}
With $w_{eco}$, $w_{env}$, and $w_{comfort}$ the weight of the economic, environmental, and comfort dimension respectively. $ref_{eco}$, $ref_{env}$, and $ref_{comfort}$ the reference value for the economic, environmental, and comfort dimension respectively.

For defining the reference value, we need to be able to compute the value of each objective function in the exact same condition. So we will need to define a basic case that will serve to evaluate the reference value for each dimension. In this case, we considered a medium comfort case with no exchanges and use of the battery storage. Which mean that the consumption profile of each member is the one defined by the devices, and when the devices have a flexibility range, we place ourselves in the middle of this range. Defined like this, we will be able to compute a Pareto front of the model by changing the weight of the different dimensions in the objective function. 

The full optimization model is a mixed integer linear programming (MILP) model, as we have binary variables for the white goods and continuous variables for the other devices. The optimization is done using the Python library Pyomo and the solver used is Gurobi.